% 连续像素坐标下的FCOS教学例；所有长度与正文算例一致。
\begin{tikzpicture}[x=1mm,y=-1mm,font=\figurefont\small,
  dist/.style={-{Stealth[length=2mm]},draw=BookBlue,line width=0.9pt}]
 \draw[draw=BookMuted,line width=0.7pt,->] (10,0)--(90,0) node[right] {$x$};
 \draw[draw=BookMuted,line width=0.7pt,->] (10,0)--(10,58) node[below] {$y$};
 \draw[draw=BookTeal,fill=BookTeal!7,line width=1pt] (20,10) rectangle (80,50);
 \draw[draw=BookMuted,dashed,line width=0.7pt] (50,10)--(50,50);
 \draw[draw=BookMuted,dashed,line width=0.7pt] (20,30)--(80,30);
 \fill[BookInk] (50,30) circle (1mm);
 \fill[BookBlue] (25,30) circle (1mm);
 \node[anchor=north west] at (51,33) {中心：$c_*=1$};
 \node[anchor=south west] at (26,27) {边缘：$c_*\approx0.302$};
 \draw[dist] (25,18)--(20,18);
 \draw[dist] (25,18)--(80,18);
 \node[anchor=south] at (22.5,17) {$5$};
 \node[anchor=south] at (57,17) {$55$};
 \draw[draw=BookBlue,line width=0.7pt] (25,18)--(25,29);
 \node[anchor=north] at (20,51) {$(20,50)$};
 \node[anchor=south] at (80,9) {$(80,10)$};
\end{tikzpicture}
